% TOP_PROBLEM: {"id": 3219, "title": "The three 4-flows conjecture", "category_id": 3, "category": {"id": 3, "name": "graph_theory", "display_name": "Graph Theory", "description": "Problems involving graphs, networks, and their properties.", "slug": "graph-theory", "order_index": 3, "created_at": "2026-07-31T15:26:25.671Z"}, "status": "open", "problem_number": "OPG-132", "set_id": 12}
% FIRST_POSTED: 2026-10-01
% ENTRY_KIND: statement_only
% UnsolvedMath ID 3219; source catalogue code OPG-132
% !TeX program = lualatex
% Definition and sources only; this file is not an LLM solution attempt.
\documentclass[11pt,a4paper]{article}
\usepackage[margin=25mm]{geometry}
\usepackage{fontspec}
\setmainfont{lmroman10-regular.otf}[BoldFont=lmroman10-bold.otf,
 ItalicFont=lmroman10-italic.otf,BoldItalicFont=lmroman10-bolditalic.otf]
\usepackage{amsmath,amssymb,mathtools}
\usepackage{unicode-math}
\setmathfont{latinmodern-math.otf}
\AtBeginDocument{\let\setminus\smallsetminus}
\usepackage{xurl,hyperref}
\hypersetup{colorlinks=true,urlcolor=blue}
\setcounter{secnumdepth}{0}
\setlength{\parindent}{0pt}
\setlength{\parskip}{5pt}
\setlength{\emergencystretch}{3em}
\begin{document}
\section*{The three 4-flows conjecture}\label{3219}
{\small UnsolvedMath ID 3219 · OPG-132 · Graph Theory · OpenGarden}
\subsection{Definitions and mathematical statement}
Let $G=(V,E)$ be a finite loopless
undirected multigraph without a
bridge. A bridge is an edge
whose removal increases the
number of connected components.
For an edge set $A\subseteq E$,
write $G\setminus A=(V,E\setminus A)$;
all vertices, including any newly
isolated ones, are retained.

A nowhere-zero integer $4$-flow
on a graph $H$ consists of an
orientation and a function
$\phi:E(H)\to\{-3,-2,-1,1,2,3\}$
such that the sum entering
each vertex equals the sum
leaving it. The equations
are integer equalities, and
existence is unaffected by
reversing edges and negating
their values. On an edgeless
graph the empty function is
allowed.

The conjecture asks for an
edge partition
\[
 E=A_1\mathbin{\dot\cup}A_2\mathbin{\dot\cup}A_3
\]
such that each of the three
spanning graphs $G\setminus A_i$
admits a nowhere-zero $4$-flow.
Parts of the partition may
be empty. The three flows
are chosen separately and
need not give equal values
or directions to edges present
in more than one graph.

Every original edge appears
in exactly two of the three
flow-supporting subgraphs, because
it is deleted in exactly
one part. Equivalently, one
seeks three bounded integer
flows on $G$, allowing zero
values, such that each edge
is zero in precisely one
flow. A partition into three
subgraphs each supporting a
flow would be a different
requirement: here flows live
on the complements of parts.
\subsection{Short English statement}
Can every bridgeless graph have its edges split into three parts so deleting any one part leaves a graph with a nowhere-zero 4-flow?
\subsection{Sources}
\begin{itemize}
\item[S1] ULAM.AI and the UnsolvedMath Contributors, supplied problems.json, record 3219 (OPG-132), OpenGarden collection. Catalogue identity and source transcription; CC BY 4.0.
\url{https://huggingface.co/datasets/ulamai/UnsolvedMath}.
\item[S2] Open Problem Garden, The three 4-flows conjecture. Original problem and discussion; the mathematical formulation below is expanded from the supplied source record.
\url{http://www.openproblemgarden.org/op/three_4_flows_conjecture}.
\end{itemize}
\end{document}
